% Part: proof-theory
% Chapter: normalization
% Section: introduction

\documentclass[../../../include/open-logic-section]{subfiles}

\begin{document}

\olfileid[id]{pt}{nor}{int}

\olsection{Pendahuluan}

Jika Anda telah !!{prove} suatu implikasi~$!A \lif !B$, Anda mungkin dapat
menggunakannya dalam !!a{proof} bagi~$!B$ dengan memakai~\Elim{\lif}.
Tentu saja, hal ini juga memerlukan !!a{proof}~$\delta_1$ bagi~$!A$.
!!^a{proof}~$\delta_1$ Anda bagi~$!A \lif !B$ kemungkinan besar berakhir
dengan~\Intro{\lif}, yang mengubah !!a{proof}~$\delta_2$ bagi~$!B$ dari
asumsi terbuka~$!A$ menjadi bukti bagi~$!A \lif !B$. Jika demikian,
!!{proof} terakhir Anda berbentuk
\[
\AxiomC{$\Discharge{!A}{x}$}
\RightLabel{$\delta_2$}
\DeduceC{$!B$}
\DischargeRule{\Intro{\lif}}{x}
\UnaryInfC{$!A \lif !B$}
\AxiomC{}
\RightLabel{$\delta_1$}
\DeduceC{$!A$}
\RightLabel{\Elim{\lif}}
\BinaryInfC{$!B$}
\DisplayProof
\]
Walaupun strategi Anda efisien karena menggunakan kembali sesuatu yang
sudah tersedia (!!a{proof} bagi~$!A \lif !B$), !!{proof} Anda sendiri tidak
efisien. Memperkenalkan $!A \lif !B$ hanya untuk segera mengeliminasinya
tampak sebagai jalan memutar. Semestinya terdapat cara yang lebih langsung
untuk !!{prove}~$!B$.

Sesungguhnya, cara seperti itu selalu ada. ``Jalan memutar'' seperti di
atas---!!a{introduction} yang langsung diikuti !!a{elimination}---selalu
dapat dihapus dari !!{proof}s deduksi alami. Proses ini disebut
\emph{normalisasi}, dan hasilnya adalah !!{proof} \emph{normal} (atau
!!a{proof} \emph{dalam bentuk normal}).

Bukti hasil ini, seperti teorema eliminasi potongan dalam kalkulus sekuen,
cukup rumit dan mengharuskan kita mempertimbangkan banyak kasus. Hasil itu
lebih sulit dibuktikan bagi \Log{N1c} klasik daripada bagi \Log{N1i}
intuisionistik. Sebagaimana eliminasi potongan menunjukkan bahwa Anda tidak
dapat !!{prove} lebih banyak dengan menggunakan aturan \Cut{} dalam kalkulus
sekuen, normalisasi menunjukkan bahwa jalan memutar tidak memungkinkan Anda
membuktikan lebih banyak daripada ketika jalan memutar itu dihindari.
Terdapat kemiripan lain: !!{proof}s normal dapat jauh lebih panjang daripada
!!{proof} terpendek dengan jalan memutar yang memberikan hasil sama,
sebagaimana !!{proof}s dalam kalkulus sekuen dapat jauh lebih pendek dengan
potongan daripada !!{proof} bebas-potongan terpendek. Sifat sub!!{formula},
yakni bahwa Anda selalu dapat !!{prove} suatu hasil hanya dengan memakai
sub!!{formula}s dari hasil itu dalam !!{proof} Anda, mengikuti normalisasi
bagi deduksi alami sebagaimana sifat tersebut mengikuti eliminasi potongan
bagi kalkulus sekuen.

\end{document}
